% ==========================================
% Archivo: sections/3_code_and_methods.tex
% ==========================================
\section{Marco Computacional}

La formulación matemática descrita requiere métodos de integración numérica explícitos para resolver el problema de valor inicial. Se han implementado dos esquemas de distinto orden de precisión: Euler Explícito ($\mathcal{O}(\Delta t)$) y Runge-Kutta clásico de 4 etapas (RK4, $\mathcal{O}(\Delta t^4)$).

El estado del sistema se discretiza en el vector de estado $\mathbf{Y}(t) = [\mathbf{r}(t), \mathbf{v}(t)]^T$. Una particularidad algorítmica fundamental incorporada en las rutinas es la detección de eventos de parada (colisión con el suelo, es decir, el cruce por cero de la componente $y$ o $z$). Para evitar errores de truncamiento severos en el instante de impacto, se implementa una interpolación lineal entre el último paso válido y el primer paso negativo.

A continuación, se presenta la implementación del método de Euler explícito, el cual avanza la solución calculando la derivada en el punto inicial del intervalo temporal:

\begin{lstlisting}[language=Python, caption={Implementación del método de Euler explícito con interpolación de impacto.}]
def euler_solver(func, t_array, y0, stop_idx=None):
    steps = len(t_array)
    dt = t_array[1] - t_array[0]
    
    y_sol = np.zeros((steps, len(y0)))
    y_sol[0] = y0
    
    for i in range(steps - 1):
        y_next = y_sol[i] + dt * func(t_array[i], y_sol[i])
        
        # Logica de parada e interpolacion lineal
        if stop_idx is not None and y_next[stop_idx] < 0:
            y_prev = y_sol[i, stop_idx]
            val_next = y_next[stop_idx]
            
            frac = y_prev / (y_prev - val_next)
            y_sol[i+1] = y_sol[i] + frac * (y_next - y_sol[i])
            return y_sol[:i+2]

        y_sol[i+1] = y_next
        
    return y_sol
\end{lstlisting}


Por su parte, la lógica del integrador \texttt{rk4\_solver} es la siguiente:
\begin{lstlisting}[language=Python, caption={Implementación del método de Runge-Kutta 4 con interpolación de impacto.}]
def rk4_solver(func, t_array, y0, stop_idx=None):
    steps = len(t_array)
    dt = t_array[1] - t_array[0]
    
    y_sol = np.zeros((steps, *np.shape(y0)))
    y_sol[0] = y0
    for i in range(steps - 1):
        t_i = t_array[i]
        y_i = y_sol[i]
        
        k1 = func(t_i, y_i)
        k2 = func(t_i + 0.5*dt, y_i + 0.5*dt*k1)
        k3 = func(t_i + 0.5*dt, y_i + 0.5*dt*k2)
        k4 = func(t_i + dt, y_i + dt*k3)
        
        y_next = y_i + (dt / 6.0) * (k1 + 2*k2 + 2*k3 + k4)
        
        # Logica de parada e interpolacion lineal
        if stop_idx is not None and y_next[stop_idx] < 0:
            y_prev_val = y_i[stop_idx]
            y_next_val = y_next[stop_idx]
            
            frac = y_prev_val / (y_prev_val - y_next_val)
            y_sol[i+1] = y_i + frac * (y_next - y_i)
            
            return y_sol[:i+2]
    
        y_sol[i+1] = y_next
    return y_sol
\end{lstlisting}

La interpolación calcula la fracción del paso de tiempo $\Delta t$ exacta en la que el cuerpo interseca $y=0$, escalando el vector completo y truncando el array de salida, optimizando así la memoria y preservando la veracidad geométrica de la trayectoria computada.


La implementación del método RK4 se utilizó exclusivamente en la modelización del lanzamiento de varios cuerpos bajo muelles viscoelásticos. Esta elección responde a la necesidad de obtener una mayor precisión en la trayectoria y en la conservación de la energía del sistema sin tener que recurrir a diferenciales de tiempo ($\Delta t$) extremadamente pequeños. A diferencia del lanzamiento de un proyectil simple, donde los métodos de primer orden resultan suficientes, la interacción entre múltiples masas acopladas por fuerzas elásticas y de amortiguamiento requiere de un algoritmo más robusto para evitar inestabilidades numéricas. El uso de RK4 permite capturar estas oscilaciones complejas de forma eficiente y físicamente coherente.
